\documentclass[12pt]{article}

\usepackage[T1]{fontenc}
\usepackage[utf8]{inputenc}
\usepackage[romanian]{babel}

\usepackage{amsmath, amssymb, amsthm}
\usepackage{mathrsfs}
\usepackage{geometry}
\geometry{a4paper, margin=2.5cm}

\title{Media (Valoarea Așteptată) a unei Variabile Aleatoare}
\author{}
\date{}

\theoremstyle{plain}
\newtheorem{definition}{Definiție}
\newtheorem{proposition}{Propoziție}
\newtheorem{example}{Exemplu}

\begin{document}

\maketitle

\section{Definiția generală a valorii așteptate}

\begin{definition}
Fie $X$ o variabilă aleatoare discretă cu valori în mulțimea cel mult numărabilă
$S = \{x_1, x_2, \dots\}$ și funcție de masă a probabilității $p(x_k) = \mathbb{P}(X = x_k)$.
\textbf{Valoarea așteptată} (sau \textbf{media}) lui $X$ este


\[
\mathbb{E}[X] := \sum_{k} x_k\, p(x_k),
\]


dacă seria este absolut convergentă, adică $\sum_k |x_k| p(x_k) < \infty$.
\end{definition}

\begin{definition}
Fie $X$ o variabilă aleatoare continuă cu densitate $f:\mathbb{R}\to[0,\infty)$.
\textbf{Valoarea așteptată} a lui $X$ este


\[
\mathbb{E}[X] := \int_{-\infty}^{\infty} x f(x)\, dx,
\]


dacă integrala Lebesgue $\int_{-\infty}^{\infty} |x| f(x)\, dx$ este finită.
\end{definition}

\section{Proprietăți fundamentale ale valorii așteptate}

\begin{proposition}[Linearitatea]
Fie $X,Y$ variabile aleatoare cu $\mathbb{E}[|X|]<\infty$, $\mathbb{E}[|Y|]<\infty$ și $a,b\in\mathbb{R}$.
Atunci


\[
\mathbb{E}[aX + bY] = a\,\mathbb{E}[X] + b\,\mathbb{E}[Y].
\]


\end{proposition}

\begin{proof}
Demonstrăm pentru cazul discret; cazul continuu este analog, cu sume înlocuite de integrale.

Presupunem că $X$ și $Y$ iau valori în mulțimile cel mult numărabile
$S_X = \{x_i\}$, $S_Y = \{y_j\}$.
Atunci


\[
\mathbb{E}[aX + bY]
= \sum_{i,j} (a x_i + b y_j)\, \mathbb{P}(X = x_i, Y = y_j).
\]


Dezvoltăm:


\[
\mathbb{E}[aX + bY]
= a \sum_{i,j} x_i\, \mathbb{P}(X = x_i, Y = y_j)
+ b \sum_{i,j} y_j\, \mathbb{P}(X = x_i, Y = y_j).
\]


Observăm că


\[
\sum_{j} \mathbb{P}(X = x_i, Y = y_j) = \mathbb{P}(X = x_i),
\]


deci


\[
\sum_{i,j} x_i\, \mathbb{P}(X = x_i, Y = y_j)
= \sum_i x_i \sum_j \mathbb{P}(X = x_i, Y = y_j)
= \sum_i x_i\, \mathbb{P}(X = x_i)
= \mathbb{E}[X].
\]


Analog,


\[
\sum_{i,j} y_j\, \mathbb{P}(X = x_i, Y = y_j)
= \sum_j y_j\, \mathbb{P}(Y = y_j)
= \mathbb{E}[Y].
\]


Prin urmare,


\[
\mathbb{E}[aX + bY] = a\,\mathbb{E}[X] + b\,\mathbb{E}[Y].
\]


\end{proof}

\begin{proposition}[Media unei constante]
Pentru orice constantă $c\in\mathbb{R}$, dacă $X \equiv c$ (adică $X(\omega)=c$ pentru orice $\omega$),
atunci


\[
\mathbb{E}[X] = c.
\]


\end{proposition}

\begin{proof}
În cazul discret, $X$ ia o singură valoare $c$ cu probabilitate $1$:


\[
\mathbb{E}[X] = \sum_x x\,\mathbb{P}(X=x) = c\cdot \mathbb{P}(X=c) = c\cdot 1 = c.
\]


În cazul continuu, densitatea este concentrată în sens de distribuție la $c$, iar integrala dă același rezultat.
\end{proof}

\begin{proposition}[Monotonicitate]
Dacă $X\le Y$ aproape sigur și $\mathbb{E}[X],\mathbb{E}[Y]$ există, atunci


\[
\mathbb{E}[X] \le \mathbb{E}[Y].
\]


\end{proposition}

\begin{proof}
Considerăm variabila $Z = Y - X \ge 0$ aproape sigur.
Atunci $\mathbb{E}[Z] \ge 0$ (deoarece integrala/suma unei funcții nenegative este nenegativă).
Dar


\[
\mathbb{E}[Z] = \mathbb{E}[Y-X] = \mathbb{E}[Y] - \mathbb{E}[X],
\]


de unde rezultă $\mathbb{E}[Y] - \mathbb{E}[X] \ge 0$, adică $\mathbb{E}[X] \le \mathbb{E}[Y]$.
\end{proof}

\section{Teorema de transport (discretă și continuă)}

\begin{proposition}[Teorema de transport – cazul discret]
Fie $X$ o variabilă aleatoare discretă cu valori în mulțimea cel mult numărabilă
$S = \{x_k\}$ și fie $g:S\to\mathbb{R}$ o funcție arbitrară.
Dacă seria $\sum_k |g(x_k)|\,\mathbb{P}(X=x_k)$ este convergentă, atunci



\[
\mathbb{E}[\,g(X)\,] = \sum_{k} g(x_k)\,\mathbb{P}(X = x_k).
\]



\end{proposition}

\begin{proof}
Variabila $g(X)$ este discretă și ia valorile $g(x_k)$ cu probabilitățile
$\mathbb{P}(X=x_k)$. Aplicăm definiția valorii așteptate pentru variabile discrete.
\end{proof}

\begin{proposition}[Teorema de transport – cazul continuu]
Fie $X$ o variabilă aleatoare continuă cu densitate $f$ și fie $g:\mathbb{R}\to\mathbb{R}$ o funcție măsurabilă.
Dacă integrala Lebesgue $\int_{\mathbb{R}} |g(x)| f(x)\,dx$ este finită, atunci



\[
\mathbb{E}[\,g(X)\,] = \int_{-\infty}^{\infty} g(x)\, f(x)\, dx.
\]



\end{proposition}

\begin{proof}
Variabila $g(X)$ are distribuție indusă de $X$ prin aplicația $g$.
Prin definiția valorii așteptate pentru variabile continue, integrăm funcția $g(x)$
ponderată cu densitatea lui $X$.
\end{proof}

\begin{example}
Dacă $g(x)=x^2$, teorema de transport dă formula generală pentru momentul de ordinul al doilea:



\[
\mathbb{E}[X^2] =
\begin{cases}
\sum_k x_k^2\,\mathbb{P}(X=x_k), & \text{(discret)} \\[6pt]
\int_{-\infty}^{\infty} x^2 f(x)\,dx, & \text{(continuu)}.
\end{cases}
\]



\end{example}

\bigskip


\section{Mediile distribuțiilor discrete fundamentale}

\subsection{Distribuția Bernoulli}

\begin{definition}
$X \sim \mathrm{Bernoulli}(p)$ dacă


\[
\mathbb{P}(X=1)=p,\quad \mathbb{P}(X=0)=1-p.
\]


\end{definition}

\begin{proposition}
Dacă $X \sim \mathrm{Bernoulli}(p)$, atunci


\[
\mathbb{E}[X] = p.
\]


\end{proposition}

\begin{proof}
Prin definiție:


\[
\mathbb{E}[X] = 0\cdot \mathbb{P}(X=0) + 1\cdot \mathbb{P}(X=1)
= 0\cdot (1-p) + 1\cdot p = p.
\]


\end{proof}

\subsection{Distribuția polinomială pe o mulțime finită}

\begin{definition}
Fie $S = \{x_1,\dots,x_k\}$ și $\mathbb{P}(X=x_i)=p_i$, cu $\sum_{i=1}^k p_i=1$.
Spunem că $X$ are distribuție polinomială pe $S$.
\end{definition}

\begin{proposition}
Media lui $X$ este


\[
\mathbb{E}[X] = \sum_{i=1}^k x_i p_i.
\]


\end{proposition}

\begin{proof}
Direct din definiția generală a valorii așteptate pentru variabile discrete.
\end{proof}

\subsection{Distribuția geometrică}

\begin{definition}
$X \sim \mathrm{Geom}(p)$ dacă


\[
\mathbb{P}(X = k) = (1-p)^{k-1} p,\quad k=1,2,\dots
\]


\end{definition}

\begin{proposition}
Dacă $X \sim \mathrm{Geom}(p)$, atunci


\[
\mathbb{E}[X] = \frac{1}{p}.
\]


\end{proposition}

\begin{proof}
Avem


\[
\mathbb{E}[X] = \sum_{k=1}^{\infty} k (1-p)^{k-1} p = p \sum_{k=1}^{\infty} k q^{k-1},
\]


unde $q = 1-p$.
Știm că pentru $|q|<1$,


\[
\sum_{k=1}^{\infty} k q^{k-1} = \frac{1}{(1-q)^2}.
\]


Aceasta se obține derivând seria geometrică $\sum_{k=0}^{\infty} q^k = \frac{1}{1-q}$.
Deci


\[
\mathbb{E}[X] = p \cdot \frac{1}{(1-q)^2} = p \cdot \frac{1}{p^2} = \frac{1}{p}.
\]


\end{proof}

\subsection{Distribuția binomială negativă}

\begin{definition}
$X \sim \mathrm{NegBin}(r,p)$ dacă


\[
\mathbb{P}(X = k) = \binom{k-1}{r-1} p^r (1-p)^{k-r},\quad k=r,r+1,\dots
\]


\end{definition}

\begin{proposition}
Dacă $X \sim \mathrm{NegBin}(r,p)$, atunci


\[
\mathbb{E}[X] = \frac{r}{p}.
\]


\end{proposition}

\begin{proof}
Interpretăm $X$ ca numărul de încercări necesare pentru a obține $r$ succese,
unde fiecare încercare este Bernoulli$(p)$ și încercările sunt independente.
Fie $X_i$ numărul de încercări necesare pentru a obține succesul $i$ după succesul $i-1$.
Atunci $X_i \sim \mathrm{Geom}(p)$, independente, și


\[
X = X_1 + \cdots + X_r.
\]


Prin linearitatea valorii așteptate și rezultatul pentru geometrică:


\[
\mathbb{E}[X] = \sum_{i=1}^r \mathbb{E}[X_i] = r \cdot \frac{1}{p} = \frac{r}{p}.
\]


\end{proof}

\subsection{Distribuția Poisson}

\begin{definition}
$X \sim \mathrm{Poisson}(\lambda)$ dacă


\[
\mathbb{P}(X = k) = e^{-\lambda} \frac{\lambda^k}{k!},\quad k=0,1,2,\dots
\]


\end{definition}

\begin{proposition}
Dacă $X \sim \mathrm{Poisson}(\lambda)$, atunci


\[
\mathbb{E}[X] = \lambda.
\]


\end{proposition}

\begin{proof}
Calculăm:


\[
\mathbb{E}[X] = \sum_{k=0}^{\infty} k e^{-\lambda} \frac{\lambda^k}{k!}
= e^{-\lambda} \sum_{k=1}^{\infty} k \frac{\lambda^k}{k!}.
\]


Pentru $k\ge1$, $k \frac{\lambda^k}{k!} = \lambda \frac{\lambda^{k-1}}{(k-1)!}$, deci


\[
\mathbb{E}[X] = e^{-\lambda} \lambda \sum_{k=1}^{\infty} \frac{\lambda^{k-1}}{(k-1)!}
= \lambda e^{-\lambda} \sum_{j=0}^{\infty} \frac{\lambda^{j}}{j!}
= \lambda e^{-\lambda} e^{\lambda} = \lambda.
\]


\end{proof}

\subsection{Distribuția hipergeometrică}

\begin{definition}
$X \sim \mathrm{Hypergeom}(N,M,n)$ dacă


\[
\mathbb{P}(X = k)
= \frac{\binom{M}{k} \binom{N-M}{n-k}}{\binom{N}{n}},
\quad k=0,1,\dots,n,
\]


unde $N$ este mărimea populației, $M$ numărul de ``succese'' în populație,
iar $n$ mărimea eșantionului extras fără înlocuire.
\end{definition}

\begin{proposition}
Dacă $X \sim \mathrm{Hypergeom}(N,M,n)$, atunci


\[
\mathbb{E}[X] = n \frac{M}{N}.
\]


\end{proposition}

\begin{proof}
Considerăm indicatorii $I_i$ pentru fiecare element extras, $i=1,\dots,n$:
$I_i = 1$ dacă al $i$-lea element extras este succes, $0$ altfel.
Atunci


\[
X = I_1 + \cdots + I_n.
\]


Prin linearitate:


\[
\mathbb{E}[X] = \sum_{i=1}^n \mathbb{E}[I_i].
\]


Dar $\mathbb{E}[I_i] = \mathbb{P}(I_i=1)$ este probabilitatea ca un element extras
să fie succes. Deoarece extragerile sunt simetrice, această probabilitate este


\[
\mathbb{P}(I_i=1) = \frac{M}{N}.
\]


Deci


\[
\mathbb{E}[X] = \sum_{i=1}^n \frac{M}{N} = n \frac{M}{N}.
\]


\end{proof}

\subsection{Distribuția multinomială}

\begin{definition}
$(X_1,\dots,X_k) \sim \mathrm{Multinomial}(n; p_1,\dots,p_k)$ dacă


\[
\mathbb{P}(X_1=x_1,\dots,X_k=x_k)
= \frac{n!}{x_1!\cdots x_k!} p_1^{x_1}\cdots p_k^{x_k},
\]


unde $x_1+\cdots+x_k = n$.
\end{definition}

\begin{proposition}
Dacă $(X_1,\dots,X_k) \sim \mathrm{Multinomial}(n; p_1,\dots,p_k)$, atunci


\[
\mathbb{E}[X_i] = n p_i,\quad i=1,\dots,k.
\]


\end{proposition}

\begin{proof}
Interpretăm $X_i$ ca numărul de apariții ale categoriei $i$ în $n$ încercări independente,
fiecare încercare având probabilitatea $p_i$ de a fi în categoria $i$.
Atunci $X_i$ are distribuție binomială $\mathrm{Bin}(n,p_i)$, iar media unei binomiale este $np_i$.
Riguros, putem scrie $X_i = \sum_{j=1}^n I_{ij}$, unde $I_{ij}$ este indicatorul că
încercarea $j$ a căzut în categoria $i$. Atunci


\[
\mathbb{E}[X_i] = \sum_{j=1}^n \mathbb{E}[I_{ij}] = \sum_{j=1}^n p_i = n p_i.
\]


\end{proof}

\section{Mediile distribuțiilor continue fundamentale}

\subsection{Distribuția uniformă pe un interval}

\begin{definition}
$X \sim \mathrm{Unif}(a,b)$ dacă


\[
f(x) =
\begin{cases}
\dfrac{1}{b-a}, & x \in [a,b], \\[6px]
0, & \text{altfel}.
\end{cases}
\]



\end{definition}

\begin{proposition}
Dacă $X \sim \mathrm{Unif}(a,b)$, atunci


\[
\mathbb{E}[X] = \frac{a+b}{2}.
\]


\end{proposition}

\begin{proof}
Prin definiție:


\[
\mathbb{E}[X] = \int_{-\infty}^{\infty} x f(x)\, dx
= \int_a^b x \frac{1}{b-a}\, dx
= \frac{1}{b-a} \int_a^b x\, dx.
\]


Calculăm integrala:


\[
\int_a^b x\, dx = \left.\frac{x^2}{2}\right|_a^b = \frac{b^2-a^2}{2}
= \frac{(b-a)(b+a)}{2}.
\]


Deci


\[
\mathbb{E}[X] = \frac{1}{b-a} \cdot \frac{(b-a)(b+a)}{2} = \frac{a+b}{2}.
\]


\end{proof}

\subsection{Distribuția exponențială}

\begin{definition}
$X \sim \mathrm{Exp}(\lambda)$ dacă


\[
f(x) =
\begin{cases}
\lambda e^{-\lambda x}, & x \ge 0,\\
0, & x < 0.
\end{cases}
\]


\end{definition}

\begin{proposition}
Dacă $X \sim \mathrm{Exp}(\lambda)$, atunci


\[
\mathbb{E}[X] = \frac{1}{\lambda}.
\]


\end{proposition}

\begin{proof}


\[
\mathbb{E}[X] = \int_0^{\infty} x \lambda e^{-\lambda x}\, dx.
\]


Facem schimbarea de variabilă $t = \lambda x$, $x = t/\lambda$, $dx = dt/\lambda$:


\[
\mathbb{E}[X] = \int_0^{\infty} \frac{t}{\lambda} \lambda e^{-t} \frac{1}{\lambda}\, dt
= \frac{1}{\lambda} \int_0^{\infty} t e^{-t}\, dt.
\]


Dar


\[
\int_0^{\infty} t e^{-t}\, dt = \Gamma(2) = 1!,
\]


deci


\[
\mathbb{E}[X] = \frac{1}{\lambda}.
\]


\end{proof}

\subsection{Distribuția Gamma}

\begin{definition}
$X \sim \mathrm{Gamma}(\alpha,\lambda)$ dacă


\[
f(x) =
\begin{cases}
\dfrac{\lambda^\alpha}{\Gamma(\alpha)} x^{\alpha - 1} e^{-\lambda x}, & x \ge 0,\\[6pt]
0, & x < 0.
\end{cases}
\]


\end{definition}

\begin{proposition}
Dacă $X \sim \mathrm{Gamma}(\alpha,\lambda)$, atunci


\[
\mathbb{E}[X] = \frac{\alpha}{\lambda}.
\]


\end{proposition}

\begin{proof}


\[
\mathbb{E}[X] = \int_0^{\infty} x \frac{\lambda^\alpha}{\Gamma(\alpha)} x^{\alpha-1} e^{-\lambda x}\, dx
= \frac{\lambda^\alpha}{\Gamma(\alpha)} \int_0^{\infty} x^{\alpha} e^{-\lambda x}\, dx.
\]


Schimbăm variabila $t = \lambda x$, $x = t/\lambda$, $dx = dt/\lambda$:


\[
\mathbb{E}[X] = \frac{\lambda^\alpha}{\Gamma(\alpha)}
\int_0^{\infty} \left(\frac{t}{\lambda}\right)^{\alpha} e^{-t} \frac{1}{\lambda}\, dt
= \frac{\lambda^\alpha}{\Gamma(\alpha)} \cdot \frac{1}{\lambda^{\alpha+1}}
\int_0^{\infty} t^{\alpha} e^{-t}\, dt.
\]


Dar


\[
\int_0^{\infty} t^{\alpha} e^{-t}\, dt = \Gamma(\alpha+1) = \alpha \Gamma(\alpha).
\]


Deci


\[
\mathbb{E}[X] = \frac{1}{\lambda} \cdot \frac{\Gamma(\alpha+1)}{\Gamma(\alpha)}
= \frac{1}{\lambda} \cdot \alpha = \frac{\alpha}{\lambda}.
\]


\end{proof}

\subsection{Distribuția normală}

\begin{definition}
$X \sim \mathcal{N}(\mu,\sigma^2)$ dacă


\[
f(x) = \frac{1}{\sigma \sqrt{2\pi}}
\exp\!\left( -\frac{(x - \mu)^2}{2\sigma^2} \right).
\]


\end{definition}

\begin{proposition}
Dacă $X \sim \mathcal{N}(\mu,\sigma^2)$, atunci


\[
\mathbb{E}[X] = \mu.
\]


\end{proposition}

\begin{proof}
Facem schimbarea de variabilă $Z = \frac{X-\mu}{\sigma}$.
Atunci $Z \sim \mathcal{N}(0,1)$, iar $X = \mu + \sigma Z$.
Prin linearitatea valorii așteptate:


\[
\mathbb{E}[X] = \mathbb{E}[\mu + \sigma Z]
= \mu + \sigma \mathbb{E}[Z].
\]


Rămâne să arătăm că $\mathbb{E}[Z]=0$ pentru $Z\sim\mathcal{N}(0,1)$.
Densitatea lui $Z$ este


\[
\varphi(z) = \frac{1}{\sqrt{2\pi}} e^{-z^2/2},
\]


care este funcție pară: $\varphi(-z)=\varphi(z)$.
Atunci


\[
\mathbb{E}[Z] = \int_{-\infty}^{\infty} z \varphi(z)\, dz.
\]


Funcția $z\varphi(z)$ este impară, deci integrala pe $\mathbb{R}$ este $0$.
Prin urmare $\mathbb{E}[Z]=0$ și


\[
\mathbb{E}[X] = \mu.
\]


\end{proof}

\section{Legături între distribuții: Exp și Gamma}

\begin{proposition}
Dacă $X_1,\dots,X_n$ sunt independente și $X_i \sim \mathrm{Exp}(\lambda)$,
atunci


\[
S_n = X_1 + \cdots + X_n \sim \mathrm{Gamma}(n,\lambda),
\]


iar


\[
\mathbb{E}[S_n] = \frac{n}{\lambda}.
\]


\end{proposition}

\begin{proof}
Fiecare $X_i$ are distribuție Gamma cu parametri $(\alpha=1,\lambda)$.
Se poate arăta (prin convoluție sau prin transformate Laplace) că suma a două variabile
Gamma independente cu același parametru de rată $\lambda$ și forme $\alpha_1,\alpha_2$
este Gamma$(\alpha_1+\alpha_2,\lambda)$.
Prin inducție, $S_n \sim \mathrm{Gamma}(n,\lambda)$.

Pentru medie, folosim linearitatea:


\[
\mathbb{E}[S_n] = \sum_{i=1}^n \mathbb{E}[X_i]
= n \cdot \frac{1}{\lambda} = \frac{n}{\lambda},
\]


în acord cu formula pentru Gamma$(n,\lambda)$.
\end{proof}

\end{document}
